\pageheading{DILATION\; /\; GROWTH THEOREM}{The exact order has positive lower type}
The single-channel result extends to finite positive feedback banks. Here a
bank has weights $h_j>0$ with $\sum_jh_j=1$, contractions
$0<q_j,r_j<1$, and normalized solution
\[
F'(z)=1+\sum_jh_jF(q_jz)F(r_jz),\qquad F(0)=0.
\]
Write $M(R)=M_F(R)$ and let $q_j^{\rho_j}+r_j^{\rho_j}=1$.
Then, with $\rho=\max_j\rho_j$,
\[
\boxed{0<\liminf_{R\to\infty}\frac{\log M(R)}{R^\rho}
\le\limsup_{R\to\infty}\frac{\log M(R)}{R^\rho}<\infty.}
\]
Thus logarithmic growth is trapped between two positive multiples of
$R^\rho$. In particular,
\[
\lim_{R\to\infty}\frac{\log\log M(R)}{\log R}=\rho.
\]
The theorem does not assert that $\log M(R)/R^\rho$ has a limit.
\fig{08_growth_type}{Finite-radius logarithmic growth and its proved lower
comparison. The normalized curves illustrate scale dependence; the plot is
not a proof of convergence to a type constant. Analytic truncation envelopes
and ordinary high-precision rounding are recorded separately.}
\claim{A lower bound at the exact exponent}
Choose a maximizing channel $(h,q,r)$ and integrate over a fixed-width
terminal interval. Let
\[
L=1/h,\quad C=\max\left(\frac{qL}{1-q},\frac{rL}{1-r}\right),\quad
B=L+\frac{2}{\min(q,r)},\quad c_-=\frac{\log2}{(B+C)^\rho}.
\]
The result is the explicit all-radius bound
\[
\log M(R)\ge c_-(R+C)^\rho,\qquad R\ge2.
\]
Unlike a fixed fractional window, the terminal window loses only an additive
radius. The shift $C$ absorbs that loss without changing the exponent.

\pageheading{DILATION\; /\; PROOF}{A translated power comparison}
Write $H(R)=\log M(R)$. Positivity gives
\[
M(R)\ge h\int_{R-L}^{R}M(qt)M(rt)\dd t
\ge M(q(R-L))M(r(R-L)),\qquad R\ge B.
\]
The choice of $C$ ensures
\[
q(R-L)+C\ge q(R+C),\qquad r(R-L)+C\ge r(R+C).
\]
On the seed interval $[2,B]$, $H(R)\ge\log R\ge\log2\ge c_-(R+C)^\rho$.
If this comparison holds at the two child radii, then
\[
\begin{aligned}
H(R)&\ge H(q(R-L))+H(r(R-L))\\
&\ge c_-\{[q(R-L)+C]^\rho+[r(R-L)+C]^\rho\}\\
&\ge c_-(q^\rho+r^\rho)(R+C)^\rho=c_-(R+C)^\rho.
\end{aligned}
\]
Let $Q_* =\max(q,r)$ and $d=(1-Q_*)B+Q_*L>0$. Each child radius is at
most $R-d$ and at least two. Induction on successive intervals of length
$d$ propagates the seed comparison to every $R\ge B$. The argument uses
monotonicity of positive powers and works also when $\rho<1$.

\claim{Explicit upper constants are available too}
Set $Q=\max_j\max(q_j,r_j)$, $m=\min_j\min(q_j,r_j)$,
$\Delta=\log(1/m)$, and $B_+=\max_j\log(1/(q_jr_j))+\log2$.
Choose an integer $K$ with $Q^{-K}\ge3/m$ and put
\[
V_K=(2^{K+1}-1)\log2+(2^{K+1}-K-2)\log(1/Q),
\]
\[
c_+=e^{-\rho}(V_K+1+\Delta+B_+).
\]
Then $\log M(R)\le c_+R^\rho-\log R-B_+$ for $R\ge e$.
The tangent majorant $M(1)<2$, followed by the integral recurrence on radii
$Q^{-k}$, bounds the initial seed by $V_K$. The logarithmic-radius comparison
then supplies the displayed upper constant. The full induction is in Appendix A.

These constants are constructive and conservative. Their purpose is to
establish two-sided growth, not to estimate the optimal type.
